\section{Experiment and metric details}
\label{app:experiment-details}

\subsection{Information-condition matrix}

% CLAIM: C03,C04,C05,C08
\Cref{tab:phase-a-methods} makes explicit which target components and future
samples enter each Phase~A condition.  ``Truth'' denotes analytic values
available only for the three analytic references.  The P row is the only
condition that supplies neither a velocity nor an acceleration target.

\begin{table}[ht]
  \centering
  \caption{Phase~A target-state and derivative information conditions.  Every
  condition uses the same one-axis limits and
  \(\mathrm{target}[k]\mapsto\mathrm{output}[k+1]\) follower clock.}
  \label{tab:phase-a-methods}
  \small
  \begin{tabular}{llll}
    \toprule
    Condition & Target components & Source & Online status \\
    \midrule
    P & \(p,0,0\) & position sample & causal \\
    PV truth & \(p,v^\star,0\) & analytic derivative & oracle \\
    PVA truth & \(p,v^\star,a^\star\) & analytic derivative & oracle \\
    PV/PVA backward & P/V or P/V/A & historical difference & causal \\
    PV/PVA centered offline & P/V or P/V/A & centered difference & noncausal \\
    PV/PVA centered causal & P/V or P/V/A & delayed centered/propagated & causal \\
    Next-cycle oracle & P/V/A at \(t_{k+1}\) & analytic future state & noncausal \\
    \bottomrule
  \end{tabular}
\end{table}

All analytic following rows use samples
\([\PhaseAEvaluationStartIndex,\PhaseAEvaluationStopIndex)\).  Derivative-quality rows
use a common interior ending two samples before the analytic sequence
boundary so that endpoint handling does not advantage one stencil.  The
development trace uses
\([\PhaseAEvaluationStartIndex,\CSVRowCount)\).  Target diagnostics align the target
consumed at \(k-1\) with the output recorded at \(k\), consistent with the
one-step interface.

\subsection{Metrics and target diagnostics}

For an evaluation sequence of length \(N\), reference position \(p_i\), and
output position \(y_i\), the unshifted metrics are
\[
  \RMSE =
  \sqrt{\frac{1}{N}\sum_{i=1}^{N}(y_i-p_i)^2},
  \qquad
  e_{\max}=\max_i |y_i-p_i|.
\]
Best grid lag searches integer shifts in the registered range and reports the
shift that minimizes the aligned RMSE.  The lag is reported in milliseconds;
positive values mean the output is late.  No statistical significance test is
applied to the three deterministic analytic reference cases.

Raw target point admissibility and target projection are measured before
following.
For PVA finite-difference rows, the acceleration diagnostic is the magnitude
of the raw differentiated target component.  It is not acceleration truth.
Output \texttt{new\_jerk}, an acceleration difference divided by the sample
period, and the maximum jerk over an exact command profile are three distinct
quantities and are not interchanged.

\subsection{Evidence lifecycle and frozen denominators}

% CLAIM: C10,C13
Phase~A aggregate CSV files and their manifests are checked-in current
evidence; their historical manifest records the generator commit, input
hashes, package versions, and a dirty generation worktree.  The v3 result layer
instead uses a frozen artifact index and byte hashes.  The audit re-streamed
the bounded locked-test tables rather than copying prose summaries.  It
preserved the \VThreeTrajectoryCount{} whole-trajectory denominator and did not delete incomplete
pairs from the formal design.

The direct-command population contains \VThreeDirectCycleCount{} recorded
command states.  Across \VThreeTrajectoryCount{} trajectories, the
adjacent-transition population is \VThreeDirectTransitionCount{} because each
trajectory contributes one fewer state pair than command states.  Runtime
discards \VThreeRuntimeWarmupSamples{} warm-up samples within each trajectory,
leaving \VThreeRuntimeCycleCount{} timed cycles.  These definitions explain
why constraint, transition, and runtime counts differ and prevent a zero count
from being reported against the wrong denominator.

% CLAIM: N03
V4 used fresh identities and a same-follower P/PV/PVA matrix without
reusing any V1, V2, or V3 identity.  Its retained observed effect is
non-confirmatory because the frozen validity and hard-runtime gates were not
all satisfied.  Another confirmation would require a new protocol, identities,
seeds, and test set; the V4 test may not be rerun.
